\documentclass[11pt]{article}
\usepackage[T1]{fontenc}
\usepackage{lmodern}
\usepackage{amsmath,amssymb,amsthm,mathtools,mathrsfs}
\usepackage[margin=1in]{geometry}
\usepackage{booktabs,array}
\usepackage[colorlinks=true,linkcolor=blue,citecolor=blue,urlcolor=blue]{hyperref}
\usepackage{enumitem}
\allowdisplaybreaks
\setlength{\emergencystretch}{2em}
\numberwithin{equation}{section}
\newtheorem{theorem}{Theorem}[section]
\newtheorem{proposition}[theorem]{Proposition}
\newtheorem{lemma}[theorem]{Lemma}
\newtheorem{corollary}[theorem]{Corollary}
\theoremstyle{definition}
\newtheorem{definition}[theorem]{Definition}
\theoremstyle{remark}
\newtheorem{remark}[theorem]{Remark}
\DeclareMathOperator{\Rea}{Re}
\DeclareMathOperator{\Ima}{Im}
\DeclareMathOperator{\Dom}{Dom}
\DeclareMathOperator{\supp}{supp}
\newcommand{\dd}{\,d}
\newcommand{\Lag}{\mathscr L}
\newcommand{\RH}{\mathrm{RH}}
\title{Constrained operator squares and stationary Laguerre energy}
\author{Drafted for Edward Baker}
\date{Draft, 14 October 2026}
\begin{document}
\maketitle
\begin{abstract}
We develop a self-contained account of constrained operator squares for the
Riemann theta heat flow and for integrated imaginary-time trial states.
The theta background includes the precise normalization of the heat family,
the de Bruijn--Newman threshold, the finite-collision reduction of the
Riemann hypothesis, and a two-level stationary-sign criterion covering
every finite multiplicity. We derive paired Dirac representations of the
physical Laguerre forms and explain why a positive generator leaves an
indefinite oscillatory observation. Local linear stationary multipliers
cannot turn that observation into a positive operator identity on the full
paired core; positive smooth source controls can also preserve finitely
many declared linear observations while both stationary signs fail.
Full imaginary Jacobi gluing remains an arithmetic constraint requiring a
quantitative signed consequence. Separately, Gaussian translation channels
supply the corrected smoothing machinery with its product defect and
normal-ordering contraction, and residual-corrected heat identities retain
the variational meaning of approximate trajectories. A positive-interaction
bound and exact finite-time diagonal identities delineate useful quantum
subclasses. The final discussion states the remaining bottlenecks and
possible constructions. Numerical work is referenced in the supporting
notes. No global theta sign theorem or proof of the Riemann hypothesis is
claimed.
\end{abstract}
\tableofcontents

\section{Purpose and notation}\label{sec:purpose}
The purpose is to collect the setup and the conditional route to the
Riemann hypothesis in the same document as investigation 20's operator
identities and obstructions. The background was previously distributed
between the parent collision reductions \cite{parentcollision}, the
stationary-sign and paired-kernel notes \cite{stationarynote,kernelnote},
the upstream manuscript \cite{upstreamtransport}, and the investigation's first research notes
\cite{diracnote,residualnote,localobstruction,smoothingnote}.
The Euler-factor control manuscript \cite{eulercontrols} uses the same
heat family, but tests which shared hypotheses fail to force real zeros;
it does not supply the missing theta sign. Here the arithmetic question
is whether the complete theta preparation forces a signed observation to
be nonnegative under its physical stationary constraint.

The motivation comes from Edward B. Baker III's continuum sigma notes
of November 2017, February 2018 and August 2018
\cite{baker2017,baker2018heat,baker2018dirac}. Their source review
\cite{seedreview} distinguishes valid heat/resolvent identities from
unfinished smoothing and square-completion steps. The useful method is
to retain the constraint and every correction when completing a square.
The theta and quantum applications have different generators, domains and
success criteria. There is no transfer theorem from a quantum variational
bound to the theta stationary signs.

Throughout the theta sections, $t$ is Newman time and $x$ the real spectral
coordinate. We write $H(t,x)=H_t(x)$ and
$H^{(j)}(t,x)=\partial_x^jH(t,x)$; a prime always differentiates $x$
at fixed $t$. The subscript on $H_t$ denotes a time slice, not a derivative.
The index $n$ labels theta-series terms, whereas $j$ labels a spatial
derivative level. Our two physical targets use $j=0,1$.
In the angular-lift section, $z$ is explicitly reassigned to the balanced
angular coordinate, and $x$ remains the spectral coordinate.
In the quantum sections $s$ is dissipative imaginary time and $T$ is a
finite terminal time. Inner products are conjugate-linear in their first
argument. Selfadjoint unbounded operators are used only on the explicitly
stated vector or form domains.

The basic logical distinction is between a positive operator on its whole
Hilbert space and a positive observation on a prescribed arithmetic state.
A multiplier that is indefinite on the whole space can have a positive
expectation on a constrained family. Conversely, the existence of a
positive square for an evolution generator does not decide an oscillatory
matrix element of that evolution. The construction problem is to display
the extra state constraint and use it quantitatively.


\section{The theta heat flow and the exact route to RH}
\label{sec:background-heat}

The scalar observation in this investigation is the de Bruijn--Newman
heat deformation of the Riemann $\xi$ function. We use $H(t,z)$ for
the joint function and $H_t(z)=H(t,z)$ for a time slice; a subscript $t$
on $H_t$ is therefore a parameter, whereas $\partial_tH$ is a derivative.
For $\Rea s>1$, $\zeta(s)=\sum_{n\ge1}n^{-s}$, and elsewhere
we use its meromorphic continuation. Its nontrivial zeros lie in the
critical strip $0<\Rea s<1$. The completed function below is entire;
its zeros are precisely those nontrivial zeros. Set
\begin{align}
 \xi(s)&=\frac{s(s-1)}2\pi^{-s/2}\Gamma(s/2)\zeta(s),\nonumber\\
 \Phi(u)&=\sum_{n\ge1}
 \bigl(2\pi^2n^4e^{9u}-3\pi n^2e^{5u}\bigr)
 e^{-\pi n^2e^{4u}},\qquad u\ge0,\label{eq:bg-Phi}\\
 H(t,z)&=\int_0^\infty e^{tu^2}\Phi(u)\cos(zu)\,du,
 &H(0,z)&=\frac18\xi\!\left(\frac12+\frac{iz}2\right).
 \label{eq:bg-H}
\end{align}
This normalization is the one in the parent collision reduction and in
the standard heat-flow literature \cite{polymath,newmanwu,parentcollision}.
The theta functional equation supplies an even extension $\Phi_e$ of
$\Phi$ to the full real line. Each term in \eqref{eq:bg-Phi} is strictly
positive on $u\ge0$, since
$2\pi n^2e^{4u}>3$. The even extension is consequently positive.

The double exponential decay of the complete theta source dominates
every fixed Gaussian, exponential, and polynomial factor on $u\ge0$.
Thus the integral and all fixed parameter derivatives converge locally
uniformly in the complex variables $t,z$. In particular $H$ is jointly
entire, real on real $(t,x)$, even in $z$, and satisfies
\begin{equation}
 \partial_t H=-\partial_z^2H.
 \label{eq:bg-backward-heat}
\end{equation}
The minus sign matters: increasing Newman time is backward heat
evolution. Moreover,
\begin{equation}
 H(t,iy)=\int_0^\infty e^{tu^2}\Phi(u)\cosh(yu)\,du>0
 \qquad(t,y\in\mathbb R),
 \label{eq:bg-imaginary-positive}
\end{equation}
so no time slice is identically zero and every spatial zero has finite
multiplicity.

\subsection{Imported global inputs and a finite collision reduction}

The following classical inputs concern the genuine theta flow; they are
not consequences of an abstract operator square. The threshold theorem
gives a finite real constant $\Lambda$ such that
\begin{equation}
 H_t\text{ has only real zeros}\quad\Longleftrightarrow\quad t\ge\Lambda.
 \label{eq:bg-threshold}
\end{equation}
De Bruijn's strip contraction maps a strip of half-width $a$ at time
$t_0$ into the strip of half-width
$\sqrt{\max\{a^2-2h,0\}}$ at time $t_0+h$, $h\ge0$.
These facts and the Riemann functional equation give
\begin{equation}
 \mathrm{RH}\quad\Longleftrightarrow\quad
 H_0\text{ has only real zeros}\quad\Longleftrightarrow\quad\Lambda\le0.
 \label{eq:bg-RH-threshold}
\end{equation}
Under $s=1/2+iz/2$, the critical line corresponds to real $z$,
and the original critical strip corresponds to $|\Ima z|<1$.
Strip contraction therefore also gives the classical bound $\Lambda\le1/2$.
Rodgers--Tao prove $\Lambda\ge0$ \cite{rodgerstao}, so RH would give
$\Lambda=0$; that lower bound is not needed in the collision argument.

We also import two more specific inputs from \cite{polymath}.
First, Theorem~1.5(i) supplies an absolute $C$ such that a zero
$H_t(x+iy)=0$ with $0<t\le1/2$ and $|x|\ge e^{C/t}$ is real.
Second, Proposition~3.1 gives local Hermite splitting: at a real zero
of multiplicity $m\ge2$ at $(T,a)$, the $m$ nearby zeros have expansions
\begin{equation}
 a+\sqrt{2(t-T)}\,\lambda_k+O(|t-T|),\qquad 1\le k\le m,
 \label{eq:bg-Hermite}
\end{equation}
where $\lambda_k$ are the distinct real roots of the probabilists'
Hermite polynomial
\[
 \operatorname{He}_m(z)=(-1)^m e^{z^2/2}
             \frac{d^m}{dz^m}e^{-z^2/2}.
\]
A nonzero $\lambda_k$
produces a nonreal root immediately before $T$.

\begin{proposition}[A positive threshold produces a finite collision]
\label{prop:bg-finite-collision}
If $\Lambda>0$, there is $a\in\mathbb R\setminus\{0\}$ with
\begin{equation}
 H(\Lambda,a)=H_x(\Lambda,a)=0.
 \label{eq:bg-positive-collision}
\end{equation}
\end{proposition}
\begin{proof}
Take $\Lambda/2\le t_k<\Lambda$ tending upward to $\Lambda$, and
choose a nonreal zero $z_k$ of each $H_{t_k}$. The original critical
strip and strip contraction bound their imaginary parts. The uniform
localization input bounds their real parts by $e^{2C/\Lambda}$, using
the classical bound $\Lambda\le1/2$. A subsequence converges to a finite
$z_*$. Continuity gives $H(\Lambda,z_*)=0$, and
\eqref{eq:bg-threshold} makes $z_*=a$ real. If this zero were simple,
the holomorphic implicit-function theorem would give a unique nearby
zero for real times near $\Lambda$. Complex conjugation would force
that zero to be real, contradicting the nearby $z_k$. Hence $H_x$ also
vanishes. Equation~\eqref{eq:bg-imaginary-positive} excludes $a=0$.
\end{proof}

The uniformity of localization for times bounded away from zero is
essential here. Merely knowing a separate finite exceptional region at
each time would not prevent the chosen zeros from escaping as
$t_k\uparrow\Lambda$.

\begin{corollary}[The exact positive-time endpoint target]
\label{cor:bg-exact-route}
RH is equivalent to the absence of a real common zero of $H_t,H_t'$
at every $t>0$. If the time-zero zeros lie in a strip of half-width $a$,
it suffices to exclude such common zeros for $0<t\le a^2/2$.
\end{corollary}
\begin{proof}
If RH holds, all zeros are real at every nonnegative time. A multiple
real zero at a positive time would contradict
\eqref{eq:bg-Hermite} at a slightly earlier nonnegative time. If RH
fails, $\Lambda>0$ and Proposition~\ref{prop:bg-finite-collision}
supplies the forbidden common zero. For the restricted time range, strip
contraction makes all zeros real at $t\ge a^2/2$, and Hermite splitting
excludes a multiple root at a later time.
\end{proof}

This is the parent manuscript's endpoint reduction
\cite{parentcollision}. A better positive bound for $\Lambda$, or a
more regular coordinate description of a collision, does not by itself
exclude \eqref{eq:bg-positive-collision}. The present task is to supply
an arithmetic sign that does exclude it.

\section{Derivative levels and the stationary Laguerre criterion}
\label{sec:background-stationary}

For spatial derivatives write $H^{(j)}=\partial_x^jH$ and define
\begin{equation}
 h_j=H^{(j)},\qquad
 \mathscr L_j(H)=(\partial_xh_j)^2-h_j\partial_x^2h_j
 =(H^{(j+1)})^2-H^{(j)}H^{(j+2)}.
 \label{eq:bg-Laguerre-level}
\end{equation}
The index $j$ denotes the derivative to which the ordinary first
Laguerre expression is applied. It is not the index of a generalized
Laguerre coefficient. Only the following two levels are needed for the
criterion:
\begin{align}
 j=0:\quad &h_0=H,
 &\mathscr L_0&=H_x^2-HH_{xx},
 &\partial_xh_0=0&\iff H_x=0,\nonumber\\
 j=1:\quad &h_1=H_x,
 &\mathscr L_1&=H_{xx}^2-H_xH_{xxx},
 &\partial_xh_1=0&\iff H_{xx}=0.
 \label{eq:bg-two-levels}
\end{align}
Every $h_j$ satisfies the same backward heat equation.

For context, the usual Laguerre inequality of a real-rooted
genus-one canonical product is immediate from its logarithmic second
derivative. If
$f(x)=C x^m e^{bx}\prod_n(1-x/\rho_n)e^{x/\rho_n}$, with real
$\rho_n$, $b,C$ and $\sum_n\rho_n^{-2}<\infty$, then away from its
zeros
\begin{equation}
 \frac{(f')^2-ff''}{f^2}
 =-\partial_x^2\log|f|
 =\frac m{x^2}+\sum_n\frac1{(x-\rho_n)^2}\ge0.
 \label{eq:bg-real-root-Laguerre}
\end{equation}
Continuity handles the zeros. Our proof target is weaker than global
Laguerre nonnegativity: it asks for the two signs only on their physical
stationary sets,
\begin{align}
 H_x=0&\quad\Longrightarrow\quad\mathscr L_0=-HH_{xx}\ge0,
 \label{eq:bg-S0}\\
 H_{xx}=0&\quad\Longrightarrow\quad\mathscr L_1=-H_xH_{xxx}\ge0.
 \label{eq:bg-S1}
\end{align}
These signs are to be proved without assuming the real-rootedness that
is the desired conclusion.

\subsection{A predecessor critical curve at every even multiplicity}

\begin{lemma}[Rescaled critical chart]
\label{lem:bg-predecessor}
Let $F_t=-F_{xx}$ be real analytic near $(T,a)$ and let $F(T,\cdot)$
have exact even multiplicity $m=2r$ at $a$. Put
$F_k=\partial_x^kF(T,a)$, so $F_m\ne0$. For sufficiently small
$h=T-t>0$ there is an analytic real curve $\chi(h)$ satisfying
\begin{align}
 \chi(h)&=a-\frac{F_{m+1}}{rF_m}h+O(h^2),
 &F_x(T-h,\chi(h))&=0,\label{eq:bg-chi}\\
 F(T-h,\chi(h))&=\frac{F_m}{r!}h^r+O(h^{r+1}),
 &F_{xx}(T-h,\chi(h))&=
 \frac{F_m}{(r-1)!}h^{r-1}+O(h^r).
 \label{eq:bg-critical-jets}
\end{align}
In particular,
\begin{equation}
 FF_{xx}(T-h,\chi(h))=
 \frac{F_m^2}{r!(r-1)!}h^{2r-1}+O(h^{2r})>0.
 \label{eq:bg-critical-positive}
\end{equation}
\end{lemma}
\begin{proof}
The heat relation and analytic Taylor expansion give
\begin{equation}
 F(T-h,a+z)=\sum_{p,q\ge0}
 \frac{F_{2p+q}}{p!q!}h^pz^q.
 \label{eq:bg-heat-Taylor}
\end{equation}
Set $z=hy$. The quotient
$K(h,y)=h^{-r}F_x(T-h,a+hy)$ extends analytically to $h=0$:
a nonzero term must have $2p+q+1\ge2r$, which implies $p+q\ge r$.
The two terms of degree $r$ are
\[
 K(0,y)=\frac{F_m}{(r-1)!}y+\frac{F_{m+1}}{r!}.
\]
Its derivative in $y$ is nonzero. The real analytic implicit-function
theorem therefore gives a real analytic $y(h)$ with $K(h,y(h))=0$.
Taking $\chi(h)=a+hy(h)$ proves \eqref{eq:bg-chi}; substitution into
\eqref{eq:bg-heat-Taylor} gives \eqref{eq:bg-critical-jets} and
\eqref{eq:bg-critical-positive}.
\end{proof}

\begin{theorem}[Two stationary signs exclude every finite multiplicity]
\label{thm:bg-stationary}
Let $H_t=-H_{xx}$ be real analytic on an open set
$\Omega\subset(0,\infty)\times\mathbb R$. Suppose every spatial
zero under consideration has finite multiplicity and
\eqref{eq:bg-S0}--\eqref{eq:bg-S1} hold throughout $\Omega$.
Then there is no $(T,a)\in\Omega$ with $H(T,a)=H_x(T,a)=0$.
\end{theorem}
\begin{proof}
For an even multiplicity $m\ge2$, apply
Lemma~\ref{lem:bg-predecessor} to $F=H$. Its predecessor branch
has $H_x=0$ and $HH_{xx}>0$, violating \eqref{eq:bg-S0}.
For an odd multiplicity $m\ge3$, apply the lemma to $F=H_x$.
This is an analytic heat solution of even multiplicity $m-1$ at the
same point. Its predecessor branch has $H_{xx}=0$ and
$H_xH_{xxx}>0$, violating \eqref{eq:bg-S1}. Openness of $\Omega$
keeps both predecessor curves in the assumed sign domain, and taking
$h<T$ keeps their times positive.
\end{proof}

The theorem and its analytic chart were developed upstream in the
stationary-sign note \cite{stationarynote}. For genuine theta, a proof
of both signs at all positive times and real heights gives RH through
Corollary~\ref{cor:bg-exact-route}. It is enough to provide a sign
domain containing every possible positive threshold collision and its
earlier critical curve. For a closed rectangle or time strip, boundary
candidates require a predecessor neighborhood; a sign just at the
collision is insufficient because the relevant products vanish there.
A bounded sign theorem gives bounded collision exclusion until its
time and spatial coverage are established. The result makes no
simplicity assertion about the time-zero zeros.

The second derivative level is indispensable to this argument. The
polynomial heat solution
\begin{equation}
 F(t,x)=x^3-6(t-T)x
 \label{eq:bg-cubic-control}
\end{equation}
has a triple zero at $(T,0)$. Before $T$ it has no real stationary
point because $F_x=3x^2+6(T-t)>0$, so the first stationary sign is
vacuous. But $F_{xx}=0$ at $x=0$ and
$F_xF_{xxx}=36(T-t)>0$, violating the second sign. In fact
$\mathscr L_0(F)=3x^4+36(t-T)^2\ge0$ everywhere: even a full first
Laguerre sign can miss this local triple collision.

\subsection{Complete paired kernels for the two levels}
\label{subsec:bg-kernels}

To connect these physical signs with operator formulas, put
\begin{equation}
 m_t(u)=e^{tu^2}\Phi_e(u),\qquad
 A_t(x)=\widehat m_t(x)=2H_t(x),\qquad
 \widehat f(x)=\int_{\mathbb R}e^{ixu}f(u)\,du.
 \label{eq:bg-full-Fourier}
\end{equation}
Define the complete paired source kernels
\begin{align}
 K_{n,t}(s)&=\int_{\mathbb R}r^{2n}m_t(s+r)m_t(s-r)\,dr,
 \nonumber\\
 J_{j,t}(s)&=\int_{\mathbb R}r^2(r^2-s^2)^j
 m_t(s+r)m_t(s-r)\,dr.
 \label{eq:bg-pair-kernels}
\end{align}
The decay of the full theta source justifies all displayed double
integrals, changes of variables, and fixed polynomial insertions,
locally uniformly in time. Symmetrizing the double Fourier expression
for \eqref{eq:bg-Laguerre-level} gives the multiplier
$\tfrac12(u-v)^2(-uv)^j$. With $u=s+r$, $v=s-r$ and absolute
Jacobian two, one obtains
\begin{equation}
 \mathscr L_j(A_t;x)=4\widehat J_{j,t}(2x),\qquad
 \mathscr L_j(H_t;x)=\widehat J_{j,t}(2x).
 \label{eq:bg-Laguerre-Fourier}
\end{equation}
In particular $J_{0,t}=K_{1,t}$ and
$J_{1,t}=K_{2,t}-s^2K_{1,t}$. All kernels are even and real.
The stationary targets are therefore precisely
\begin{equation}
 \widehat J_{0,t}(2x)\ge0\ \text{when }A_t'(x)=0,
 \qquad
 \widehat J_{1,t}(2x)\ge0\ \text{when }A_t''(x)=0.
 \label{eq:bg-pair-targets}
\end{equation}
The source product in \eqref{eq:bg-pair-kernels} contains the complete
theta sum in both factors, including every cross contribution. Dropping
cross terms changes the target.

Global nonnegativity of $\widehat J_{j,t}$ would be a stronger sufficient
condition. For these real even Schwartz kernels, that condition is
equivalent to positive definiteness as a translation kernel:
$\sum_{a,b}c_a\overline{c_b}J_{j,t}(s_a-s_b)\ge0$ for all finite
choices. It is distinct from pointwise positivity of $J_{j,t}$.
For instance, $J_{0,t}>0$ follows directly from the positive source but
does not establish the sign of its cosine transform; the weight
$r^2-s^2$ already permits a signed $J_{1,t}$. The complete derivation,
the distinction from generalized Laguerre coefficients, and the local
triple limitation are recorded in \cite{kernelnote}. These formulas
specify the sign that a constrained square must produce; they do not
prove either theta target.

\section{Stationary Laguerre energy and the local observation}\label{sec:critical-energy}
The stationary conditions have a direct dynamical meaning for the
unnormalized heat solution. They concern the critical value of a spatial
slice rather than the energy integrated over that slice.

\begin{proposition}[Derivative of a squared critical value]\label{prop:critical-energy}
Let $h$ be a real $C^{1,4}$ solution of $h_t=-h_{xx}$ on a neighborhood
of a differentiable curve $x_s(t)$ with $h_x(t,x_s(t))=0$. Put
$g(t)=h(t,x_s(t))$. Then
\begin{equation}\label{eq:critical-energy}
 g'=-h_{xx},\qquad
 h_x^2-hh_{xx}=gg'=\frac12\frac{d}{dt}g^2
 \quad\text{on the curve}.
\end{equation}
If $h_{xx}\ne0$ along the curve and the displayed higher derivatives
exist, then
\begin{equation}\label{eq:critical-higher-jets}
 x_s'=\frac{h_{xxx}}{h_{xx}},\qquad
 g''=h_{xxxx}-\frac{h_{xxx}^2}{h_{xx}}.
\end{equation}
\end{proposition}
\begin{proof}
The chain rule gives $g'=h_t+x_s'h_x=-h_{xx}$. Differentiating
$h_x(t,x_s(t))=0$ gives $-h_{xxx}+x_s'h_{xx}=0$.
A further differentiation of $g'=-h_{xx}$ gives
$g''=h_{xxxx}-x_s'h_{xxx}$. These are precisely the displayed identities.
\end{proof}
Apply this proposition to $h=h_j=\partial_x^jH$, $j=0,1$.
Thus the two stationary signs require the squared critical values of
$H$ and of $H_x$ to be nondecreasing in Newman time. Equation
\eqref{eq:critical-higher-jets} supplies a possible comparison equation,
but neither term in $g''$ has an automatic favorable sign. Its denominator
also requires a separate chart at a degenerate stationary point. The
predecessor argument of the preceding section treats finite multiplicity
without assuming this ordinary chart remains nonsingular.

\subsection{The curvature term lost by spatial integration}
For any real heat solution $h$,
\begin{equation}\label{eq:local-square-defect}
 (\partial_t+\partial_x^2)h^2=2h_x^2,\qquad
 h_x^2-hh_{xx}=2h_x^2-\frac12\partial_x^2(h^2).
\end{equation}
With sufficient decay and Sobolev regularity for integration by parts,
\begin{equation}\label{eq:integrated-energy}
 \frac{d}{dt}\|h\|_{L^2(\mathbb R)}^2
 =2\|h_x\|_{L^2(\mathbb R)}^2\ge0.
\end{equation}
On $h_x=0$, however, the local Laguerre form equals
$-\partial_x^2(h^2)/2$. Spatial integration removes exactly this curvature
term. Consequently the positive integrated square cannot supply the
stationary sign without an additional local constraint.

An exact elementary control illustrates the distinction:
\begin{equation}\label{eq:elementary-critical-control}
 h(t,x)=1+ce^t\cos x,\qquad 0<c<1.
\end{equation}
It satisfies backward heat and has a positive even discrete spectral
measure. At $x=\pi$, with $C=ce^t$,
\begin{equation}
 h_x=0,\quad h=1-C,\quad h_{xx}=C,\quad
 \mathscr L_0=-C(1-C)<0\quad\text{when }0<C<1.
\end{equation}
At $C=1$ it has an ordinary double zero. The finite spectral representation
\begin{equation}
 \mathsf D=\operatorname{diag}(0,1,-1),\quad
 v=(1,\sqrt{c/2},\sqrt{c/2}),\quad
 h(t,x)=\langle v,e^{t\mathsf D^2}e^{ix\mathsf D}v\rangle
\end{equation}
retains a positive square generator and the wrong stationary sign.
This control screens the generic operator inference; it is not claimed
to obey the genuine theta arithmetic or its source decay. The smooth
Euler-factor controls provide a stronger source comparison
\cite{eulercontrols,stationarycontrols}.

\subsection{Time normalization changes the energy identity}
Let $m_0(t)=H_t(0)>0$ and $C_t(x)=H_t(x)/m_0(t)$. Put
\begin{equation}
 \mu_2(t)=\frac{\int_0^\infty u^2e^{tu^2}\Phi(u)\dd u}{m_0(t)}.
\end{equation}
Then $m_0'/m_0=\mu_2$, so
\begin{equation}\label{eq:normalized-newman}
 \partial_t C=-C_{xx}-\mu_2 C.
\end{equation}
Writing $c_j=\partial_x^jC$, a stationary curve of $c_j$ satisfies
\begin{equation}\label{eq:normalized-critical-energy}
 \frac12\frac{d}{dt}c_j(t,x_s(t))^2
 =\mathscr L_j(C;t,x_s(t))-\mu_2(t)c_j(t,x_s(t))^2.
\end{equation}
Because this normalizer depends only on time, the physical stationary
sets are unchanged. The critical-value energy evolution is changed.
A spatially varying normalizer would change the stationary sets as well.
The source normalization must therefore be fixed before stating a
monotonicity target.

\subsection{A weaker sufficient bound on stationary loss}
Nonnegative stationary Laguerre forms are a sufficient target, not a
necessary condition for excluding collisions. The upstream bounded-loss
criterion \cite{boundedloss} gives a useful alternative that keeps its
hypotheses free of division by a possibly vanishing physical observation.

\begin{proposition}[Locally bounded loss excludes collisions]\label{prop:bounded-loss}
Let $H$ be a real analytic solution of backward heat on an open
positive-time set $\Omega$, with finite spatial multiplicity at every
zero under consideration, and put
$\mathcal E=H^2+H_x^2$. Suppose nonnegative functions $C_0,C_1$ are
locally bounded on $\Omega$ and
\begin{align}
 H_x=0&\ \Longrightarrow\ \mathscr L_0(H)\ge-C_0\mathcal E,\nonumber\\
 H_{xx}=0&\ \Longrightarrow\ \mathscr L_1(H)\ge-C_1\mathcal E.
 \label{eq:bounded-loss-target}
\end{align}
Then $H=H_x=0$ has no solution in $\Omega$.
\end{proposition}
\begin{proof}
At an even multiple zero of order $2r$, the predecessor curve from the
preceding section has, for $h=T-t\downarrow0$,
$H=ch^r+O(h^{r+1})$, $H_x=0$,
$H_{xx}=rc h^{r-1}+O(h^r)$, with $c\ne0$. Hence
$HH_{xx}/\mathcal E=r/h+O(1)$, contradicting local boundedness of $C_0$.
At an odd multiple zero of order $2r+1$, apply that lemma to $H_x$.
The curve has $H_{xx}=0$, $H_x=ch^r+O(h^{r+1})$ and
$H_{xxx}=rc h^{r-1}+O(h^r)$. The heat Taylor expansion along its
$x=a+hy(h)$ chart gives $H=O(h^{r+1})$: terms with $p+q\le r$ have
$2p+q\le2r$ and thus vanish at the putative zero. Consequently
$H_xH_{xxx}/\mathcal E=r/h+O(1)$, contradicting local boundedness of $C_1$.
All divisions in this argument occur only on predecessor branches with
$\mathcal E>0$. Openness retains those branches in $\Omega$.
\end{proof}
The loss coefficients must remain locally finite through every candidate
joint zero. Defining them by an uncontrolled inverse power of
$\mathcal E$ does not establish this hypothesis. The criterion permits
some negative stationary values, but the allowance still vanishes
quadratically in the joint first jet. No such full-source theta bound
is established here. The finite-observation wrong-sign controls below
screen the stronger nonnegativity target; their negative values alone
do not refute \eqref{eq:bounded-loss-target}.


\section{Dirac squares and the physical stationary energy}
\label{sec:theta-dirac}

The useful principle in the Dirac source notes is to construct a positive
quadratic form from a first-order operator while retaining every term in
the expression to be factored.  For a selfadjoint geometric Dirac operator
$D$ with $D^2=-\Delta$, for example,
\[
 \langle f,(-\Delta)^{-1}f\rangle=\|D^{-1}f\|^2
\]
on a domain where both inverses exist.  This principle does not identify
the Newman observation with an inverse-Laplacian form.  It also does not
identify $e^{sD}$ with the positive square root $e^{-sD^2/2}$ of the
dissipative heat operator.  Newman evolution has the growing multiplier
$e^{tD^2}$ in the spectral representation; for positive $t$ it is
unbounded on an unrestricted $L^2$ space.  The theta source supplies the
analytic vectors needed here.  These distinctions, and the limitations
of the unfinished many-body square completion in the source notes,
are recorded in \cite{seedreview,diracnote}.

\subsection{An exact paired Hilbert-space formula}
\label{sec:theta-paired}

Use the genuine even source, and let
\[
 d\mu(u)=\frac12\Phi_e(u)\,du,\qquad
 \mathcal H=L^2(\mathbb R,d\mu).
\]
Multiplication by $u$ defines a selfadjoint operator $\mathsf D$ on its
maximal multiplication domain.  Set
\[
 w_t(u)=e^{tu^2/2},\quad W_t=w_t\otimes w_t,\quad
 \mathsf A=\mathsf D\otimes I-I\otimes\mathsf D,\quad
 \mathsf S=\mathsf D\otimes I+I\otimes\mathsf D.
\]
These are spectral multiplication operators; no extra Clifford
structure is needed for this representation.  Theta decay places
$W_t$ in every polynomial form domain used below, locally uniformly in
real time.  The physical observation is exactly
\begin{equation}
 H_t(x)=\langle w_t,e^{ix\mathsf D}w_t\rangle
       =\int_{\mathbb R}e^{tu^2}e^{ixu}\,d\mu(u).
 \label{eq:theta-one-body}
\end{equation}
Differentiating and symmetrizing the two spectral variables gives
\begin{align}
 \mathscr L_j(t,x)
 &=\frac12\iint_{\mathbb R^2}
     (-uv)^j(u-v)^2e^{t(u^2+v^2)}
                  \cos\bigl(x(u+v)\bigr)\,d\mu(u)d\mu(v)
 \notag\\
 &=\frac12\langle\mathsf A W_t,
       (-\mathsf D\otimes\mathsf D)^j
                       \cos(x\mathsf S)\mathsf A W_t\rangle,
 \qquad j=0,1.
 \label{eq:theta-paired-identity}
\end{align}
For completeness, the unsymmetrized difference of derivative products
has multiplier
$(-uv)^j(v^2-uv)e^{ix(u+v)}$.  Averaging with the multiplier obtained by
interchanging $u$ and $v$ produces $(u-v)^2/2$; evenness removes the
imaginary part.  This proves \eqref{eq:theta-paired-identity} with its
normalization and its extra signed factor at $j=1$.

For $j=0$, define by spectral calculus
\[
 \mathsf C_\pm=\frac{|\cos(x\mathsf S)|\pm\cos(x\mathsf S)}2.
\]
Then the complete signed expression is
\begin{equation}
 2\mathscr L_0
 =\|\mathsf C_+^{1/2}\mathsf A W_t\|^2
    -\|\mathsf C_-^{1/2}\mathsf A W_t\|^2.
 \label{eq:theta-cosine-sectors}
\end{equation}
There already is a squared first-order factor.  The unresolved input is
domination of its negative cosine sector under the physical stationary
constraint.  For $j=1$, the same positive/negative-part decomposition
can be applied to the multiplier $-uv\cos(x(u+v))$, but controlling
the cosine alone would leave the sign of $-uv$ untreated.  The complete
paired kernels and the derivative-level distinction are also discussed
in \cite{kernelnote,diracnote}.

\section{What a local stationary square cannot do}
\label{sec:theta-local-obstruction}

The following result concerns an operator identity on a full paired
core.  That requirement is stronger than an identity after expectation
in the physical vector $W_t$, and the distinction is essential.

Let $\mathcal K=L^2(\mathbb R^2,d\mu\otimes d\mu)$, with common
compact-support core $C_c^\infty(\mathbb R^2)$, and define multiplication
kernels
\begin{align*}
 k_{0,x}(u,v)&=\frac12(u-v)^2\cos\bigl(x(u+v)\bigr),\\
 k_{1,x}(u,v)&=-\frac12uv(u-v)^2\cos\bigl(x(u+v)\bigr).
\end{align*}
For a real locally essentially bounded function $m_x$, with enough
theta-weighted integrability for the displayed scalar product, put
\[
 M_{t,x}=\int e^{tu^2}m_x(u)\,d\mu(u).
\]
The stationary factors at the two derivative levels are
\[
 b_{0,x}(u)=u\sin(xu),\qquad
 b_{1,x}(u)=u^2\cos(xu),
\]
since $H_t'=-\int e^{tu^2}b_{0,x}\,d\mu$ and
$H_t''=-\int e^{tu^2}b_{1,x}\,d\mu$.  Thus
$H_t^{(j+1)}M_{t,x}$ is the expectation in $W_t$ of the multiplication
kernel
\begin{equation}
 p_{j,x}(u,v)=-\frac12\bigl(
       b_{j,x}(u)m_x(v)+m_x(u)b_{j,x}(v)\bigr).
 \label{eq:theta-stationary-kernel}
\end{equation}

\begin{theorem}[Full-core stationary-multiplier obstruction]
\label{thm:theta-local-no-go}
For each $x\ne0$ and each locally essentially bounded $m_x$, both
$K_{j,x}-P_{j,x}$, $j=0,1$, have strictly negative quadratic forms on
some vector in $C_c^\infty(\mathbb R^2)$.  In particular, on that core
neither can equal $Q^*Q$ or a sum of positive operator forms.  This
exclusion applies whether those positive operators are local or
nonlocal.
\end{theorem}

\begin{proof}
For $j=0$, put $a=\pi/|x|$.  At $(u,v)=(0,a)$ both $b_{0,x}$ factors
vanish, while $k_{0,x}(0,a)=-a^2/2$.  Pointwise evaluation of $m_x$
is unnecessary.  Choose $r_0\le a/8$ and a common essential bound $C$
for $|m_x|$ on the neighborhoods of $0$ and $a$ of radius $r_0$.
For
\[
 0<r\le\min\left\{r_0,\frac{a}{32(1+C)|x|}\right\},\qquad
 \mathcal R=[-r,r]\times[a-r,a+r],
\]
one has
\[
 |u-v|\ge3a/4,\quad \cos(x(u+v))\le-1/2,\quad
 k_{0,x}<-a^2/8.
\]
Also $|b_{0,x}(u)|\le |x|r^2$ and
$|b_{0,x}(v)|\le(9a/8)|x|r$, so that
$|p_{0,x}|\le2C|x|ar\le a^2/16$.  Consequently
\begin{equation}
 k_{0,x}-p_{0,x}\le-a^2/16\quad\hbox{a.e. on }\mathcal R.
 \label{eq:theta-negative-rectangle}
\end{equation}
The strictly positive theta measure gives positive mass to every open
rectangle.  A nonzero smooth vector supported in the interior of
$\mathcal R$ therefore has strictly negative form.

For $j=1$, put $a=\pi/(2|x|)$.  At $(a,3a)$ both $b_{1,x}$ factors
vanish, and
\[
 k_{1,x}(a,3a)=-6a^4.
\]
By continuity of $k_{1,x}$ and $b_{1,x}$ and local essential
boundedness of $m_x$, a sufficiently small open rectangle about that
point satisfies $k_{1,x}-p_{1,x}<0$ almost everywhere.  The same
compact-support argument applies.  Finally every $Q^*Q$ has
nonnegative form on its domain, so neither negative form can equal
such a positive form, independently of the locality of $Q$.
\end{proof}

This theorem includes any finite first-order multiplication ansatz
\[
 q_{\ell,x}(u,v)=\alpha_{\ell,x}(u,v)(u-v)
       +\beta_{\ell,x}(u,v)(u+v)+\gamma_{\ell,x}(u,v),
\]
with defined domains, because $\sum_\ell|q_{\ell,x}|^2$ is
nonnegative.  It also includes coefficients chosen from the genuine
theta source.  The relevant hypothesis is local boundedness of the
linear stationary multiplier near the displayed nodes.  Poles there
would require a different domain and a new boundary analysis.

There is a quantitative consequence for any multiplication identity
with an explicit remainder,
\[
 k_{0,x}=\sum_\ell|q_{\ell,x}|^2+p_{0,x}+e_x.
\]
Writing $(e_x)_-=\max\{-e_x,0\}$,
\begin{equation}
 \int (e_x)_-|W_t|^2\,d\mu\,d\mu
 \ge\frac{a^2}{16}\int_{\mathcal R}
          e^{t(u^2+v^2)}\,d\mu(u)d\mu(v)>0.
 \label{eq:theta-negative-payment}
\end{equation}
This lower bound concerns the negative part.  It does not determine the
sign of the total remainder, which may have positive contributions
elsewhere.  It identifies a signed payment that such a completion must
retain \cite{localobstruction}.

At $x=0$, $k_{0,0}=(u-v)^2/2$ is already a positive square, so the
$j=0$ obstruction is a nonzero-frequency one.  More generally, a
gluing defect may have zero expectation on the genuine theta state
while being a nonzero operator on the full core.  Removing it as an
operator would be invalid.  An identity restricted to an explicitly
arithmetic state space, or a scalar identity keeping such a gluing
term, remains outside Theorem~\ref{thm:theta-local-no-go}.  Merely
making $Q$ nonlocal while preserving the same full-core equality
$K_{j,x}-P_{j,x}=Q^*Q$ does not avoid the theorem.

\section{Finite source observations do not enforce the stationary signs}
\label{sec:theta-finite-observations}

A second obstruction concerns the information used to restrict the
source, rather than the operator domain.  It permits any prescribed
finite collection of linear source observations and local source jets.
It does not test full imaginary Jacobi gluing.

For an even source $\psi$, write
\[
 H_{\psi,t_0}(x)=\frac12\int_{\mathbb R}
             e^{t_0u^2}\psi(u)e^{ixu}\,du.
\]
Declare finitely many real linear observations
\begin{equation}
 \ell_i(\psi)=\int_{\mathbb R}K_i(u)\psi(u)\,du,
 \qquad i=1,\ldots,r,
 \label{eq:theta-finite-measurements}
\end{equation}
where $K_i$ is locally integrable and every displayed integral exists
for the base source and the perturbations below.  A complex observation
can be replaced by its two real parts.  Also declare finitely many jets
$\psi^{(k)}(s_j)$.

\begin{theorem}[Two wrong stationary signs with identical finite data]
\label{thm:theta-finite-controls}
Fix a real time $t_0$ and the finite data just described.  For every
sufficiently large $X$ there is a positive even smooth source $\psi_X$
such that $\psi_X/\Phi_e-1$ has support in a common compact set,
\[
 \ell_i(\psi_X)=\ell_i(\Phi_e),\qquad
 \psi_X^{(k)}(s_j)=\Phi_e^{(k)}(s_j),
\]
and $\psi_X/\Phi_e-1$ tends to zero in every fixed derivative supremum
norm. Every Gaussian-weighted polynomial moment remains finite. With
$\varepsilon_X=e^{-\sqrt X}$, the same source can be chosen to satisfy
\begin{align}
 H_{\psi_X,t_0}(X)&=\varepsilon_X,&
 H_{\psi_X,t_0}'(X)&=0,&
 H_{\psi_X,t_0}''(X)&=\varepsilon_X,
 \label{eq:theta-j0-control}\\
 H_{\psi_X,t_0}'(2X)&=\varepsilon_X,&
 H_{\psi_X,t_0}''(2X)&=0,&
 H_{\psi_X,t_0}'''(2X)&=\varepsilon_X.
 \label{eq:theta-j1-control}
\end{align}
Thus $\mathscr L_0(X)=\mathscr L_1(2X)=-\varepsilon_X^2<0$.
Both stationary points are nondegenerate and their wrong signs persist
locally in heat time.
\end{theorem}

\begin{proof}
Choose a symmetric bounded open set $\Omega$ avoiding neighborhoods of
the declared jet points and their reflections.  Let $V$ be the even
$C_c^\infty(\Omega)$ functions.  Restrict the observation functionals
to $q\mapsto\ell_i(\Phi_eq)$ on $V$, retain an independent basis of
their range, and choose $b_1,\ldots,b_s\in V$ dual to that basis:
\[
 \ell_i(\Phi_eb_j)=\delta_{ij},\qquad i,j=1,\ldots,s.
\]
The observations discarded as redundant on $V$ will be preserved
whenever the retained ones are preserved.

Choose a nonzero nonnegative $\phi\in V$ whose support contains an
interval.  Put
\[
 F(x)=H_{\Phi_e\phi,t_0}(x),\qquad
 I_{2k}=\frac12\int e^{t_0u^2}\Phi_e(u)\phi(u)u^{2k}\,du.
\]
Strict Cauchy--Schwarz for this positive interval-supported measure
gives $I_0I_4>I_2^2$ and $I_2I_6>I_4^2$.  For sufficiently small
fixed $d>0$,
\begin{align}
 F'(d)&=-I_2d+O(d^3)\ne0,\notag\\
 F(d)F''(0)-F(0)F''(d)
     &=\frac12(I_2^2-I_0I_4)d^2+O(d^4)\ne0,
 \label{eq:theta-first-moment-determinant}\\
 F'(d)F'''(2d)-F'(2d)F'''(d)
     &=(I_2I_6-I_4^2)d^4+O(d^6)>0.
 \label{eq:theta-second-moment-determinant}
\end{align}
For a central frequency $Y$ and fixed offset $a$, project the carrier
off the finite observation space:
\begin{equation}
 q_{a,Y}(u)=\phi(u)\cos((Y+a)u)
 -\sum_{i=1}^s
       \ell_i\bigl(\Phi_e\phi\cos((Y+a)\cdot)\bigr)b_i(u).
 \label{eq:theta-projected-carrier}
\end{equation}
It is exactly observation-null and vanishes near all jet nodes.  The
coefficients of the subtracted functions tend to zero by the
Riemann--Lebesgue lemma on the common compact support.

Take three carriers centered at $Y=X$ with offsets $-d,0,d$.  Their
responses in $(H(X),H'(X),H''(X))$ tend to
\begin{equation}
 R_0=\frac12
 \begin{pmatrix}
 F(d)&F(0)&F(d)\\
 F'(d)&0&-F'(d)\\
 F''(d)&F''(0)&F''(d)
 \end{pmatrix}.
 \label{eq:theta-three-carrier-j0}
\end{equation}
For example, the unprojected carrier response in $H(X)$ is
$\{F(a)+F(2X+a)\}/2$; the high-frequency term and the projected basis
responses tend to zero.  Differentiation supplies the other rows,
including their signs.  By
\eqref{eq:theta-first-moment-determinant}, $R_0$ is invertible.

Use another three carriers centered at $Y=2X$ with offsets $0,d,2d$.
Their responses in $(H'(2X),H''(2X),H'''(2X))$ tend to
\begin{equation}
 R_1=\frac12
 \begin{pmatrix}
 0&-F'(d)&-F'(2d)\\
 F''(0)&F''(d)&F''(2d)\\
 0&-F'''(d)&-F'''(2d)
 \end{pmatrix}.
 \label{eq:theta-three-carrier-j1}
\end{equation}
Here $F''(0)=-I_2<0$, and
\eqref{eq:theta-second-moment-determinant} makes $R_1$ invertible.

Stack the six desired derivative observations and the six carriers
into a response matrix $R_X$.  Cross-block responses tend to zero
because the two central frequencies separate by $X$.  Hence
$R_X\to\operatorname{diag}(R_0,R_1)$ and its inverses are uniformly
bounded for all sufficiently large $X$.  Let $J_X$ be the stacked
six-jet vector of the genuine heat transform and solve
\begin{equation}
 c_X=R_X^{-1}
 \left[
 (\varepsilon_X,0,\varepsilon_X,
                  \varepsilon_X,0,\varepsilon_X)^T-J_X
 \right],\qquad
 \psi_X=\Phi_e\left(1+\sum_{\nu=1}^6(c_X)_\nu q_{\nu,X}\right),
 \label{eq:theta-six-carrier-solution}
\end{equation}
where $q_{\nu,X}$ lists the six projected carriers.  The heated theta
source and its polynomial insertions are Schwartz, so $J_X$ and
$\varepsilon_X$ are $O_N(X^{-N})$ for every fixed $N$.  The coefficient
vector $c_X$ has that same decay.  The carriers are uniformly bounded,
and every fixed-order derivative grows at most polynomially in $X$.
Thus the relative perturbation and each fixed derivative tend to zero
superalgebraically.  Its supremum is eventually less than $1/2$,
giving strict positivity.  Outside the common compact support the
source is exactly $\Phi_e$, so its Gaussian moments remain finite.
The projected carriers and the solved system prove all equalities and
both negative signs.  Finally $H''(X)\ne0$ and $H'''(2X)\ne0$ permit
the two stationary curves by the implicit-function theorem; continuity
preserves the strict wrong signs for a local time interval.
\end{proof}

No effective lower bound on $X$ or uniform time radius is asserted.
The proof gives a source-class obstruction, not a counterexample for
the genuine theta preparation.  It also does not say that every
finite symbolic formula is inadequate.  A finite formula imposing an
entire functional identity, a nonlinear restriction on the whole
source class, or a statement restricted to a fixed compact frequency
region need not fall under this theorem.  A finite angular test is
covered only after its evaluated source-to-lift map has been shown to
be an observation of the form
\eqref{eq:theta-finite-measurements} or a declared local jet.
The detailed construction and a conditional analytic-source variant
are in \cite{localobstruction}.
In particular, a strict wrong stationary sign does not by itself
refute a weaker bounded-loss estimate of the form
$\mathscr L_j\ge-C_j(H^2+H_x^2)$ with a locally bounded $C_j$.
The finite-observation conclusion here excludes a zero-defect,
nonnegative stationary target based only on the declared data; it
does not exclude that weaker collision criterion.

\section{Full Jacobi gluing and the insertion that must be transported}
\label{sec:theta-gluing}

The preceding controls leave an arithmetic constraint available: the
full imaginary quasiperiod law.  To see its exact content, consider a
normally convergent coefficient preparation
\[
 \Theta_c(r,y)=\sum_{n\in\mathbb Z}
                 c_ne^{-\pi rn^2+2\pi iny}.
\]
The classical theta law in this convention is
\begin{equation}
 \Theta(r,y+ir)=e^{\pi r-2\pi iy}\Theta(r,y).
 \label{eq:theta-imaginary-law}
\end{equation}
It is a rescaling of Jacobi quasiperiodicity \cite{dlmf}.
Reindexing the series gives its defect exactly:
\begin{align}
 E_c(r,y)
 &:=e^{-\pi r}\bigl[
      \Theta_c(r,y+ir)-e^{\pi r-2\pi iy}\Theta_c(r,y)\bigr]
 \notag\\
 &=\sum_n(c_n-c_{n+1})
                  e^{-\pi r(n+1)^2}e^{2\pi iny}.
 \label{eq:theta-coefficient-defect}
\end{align}
Fourier uniqueness therefore forces $c_n=c_{n+1}$ for every $n$ when
the complete defect vanishes.  With $c_0=1$ this selects genuine theta.
Parseval gives the quantitative formula
\begin{equation}
 \int_0^1|E_c(r,y)|^2\,dy
 =\sum_n|c_n-c_{n+1}|^2e^{-2\pi r(n+1)^2}.
 \label{eq:theta-weighted-gluing-norm}
\end{equation}
At fixed $r$, the weights weaken at high index.  Exact rigidity and a
uniform estimate from a small defect are consequently different
claims.  A quantitative argument must control these weights, any
cutoff, and angular continuation.

\subsection{Balanced lift and moving-period connection}
\label{sec:theta-moving-period}

Introduce a balanced angular coordinate $z$, distinct from the real
spectral variable $x$, and define
\begin{equation}
 F(u,z)=e^{u+\pi z^2/2}\Theta(e^{4u},e^{2u}z),\qquad
 M(u,z)=\frac{(\partial_u^2-1)F(u,z)}{16},
 \label{eq:theta-balanced-insertion}
\end{equation}
where $\partial_u$ holds $z$ fixed.  Jacobi inversion gives
$F(u,z)=F(-u,-iz)$ and the same reflection for $M$.  The zero-mode
carrier is annihilated by $\partial_u^2-1$, and
$M(u,0)=\Phi_e(u)$.  The resulting physical lift
\[
 K_t(x,z)=\frac12\int_{\mathbb R}e^{tu^2+ixu}M(u,z)\,du
\]
is entire in its complex parameters and has $K_t(x,0)=H_t(x)$.
For bounded $z$, the nonzero theta modes and any fixed derivatives
are dominated at large positive $u$ by a polynomial factor times
$\exp(-\pi n^2e^{4u}/2)$; reflection pays the other tail.

Both real and imaginary balanced periods have the shape
\[
 F(u,z+p)=e^qF(u,z),\qquad
 q=\kappa(pz+p^2/2),
\]
with
\[
 \begin{array}{c|c|c|c}
 &p(u)&\kappa&p'(u)\\ \hline
 \text{real period}&\pm e^{-2u}&\pi&-2p\\
 \text{imaginary quasiperiod}&\pm i e^{2u}&-\pi&2p
 \end{array}
 \qquad p''=4p.
\]
Set $B_\kappa=\partial_z-\kappa z$.  Transport through the physical
insertion yields the exact identity
\begin{equation}
 e^{-q}M(u,z+p)
 =M(u,z)+\frac{-2p'B_\kappa F_u-p''B_\kappa F
                          +(p')^2B_\kappa^2F}{16}.
 \label{eq:theta-inserted-law}
\end{equation}
Indeed, if $T_pf(u,z)=f(u,z+p(u))$ and
$D_u=\partial_u|_z$, then
\[
 T_pD_u=(D_u-p'\partial_z)T_p,\qquad
 e^{-q}(D_u-p'\partial_z)e^q=D_u-p'B_\kappa.
\]
Expanding $\{(D_u-p'B_\kappa)^2-1\}F/16$ proves
\eqref{eq:theta-inserted-law}.  In particular the connection term
$-p''B_\kappa F$ and the scalar term $-\kappa$ in
\[
 B_\kappa^2=\partial_z^2-2\kappa z\partial_z
                      +\kappa^2z^2-\kappa
\]
are required.  Holding the period fixed during differentiation loses
physical terms.

Symmetrizing $p$ and $-p$ cancels exactly the odd connection terms:
\begin{equation}
 \frac12\sum_{\sigma=\pm1}
    e^{-\kappa(\sigma pz+p^2/2)}M(u,z+\sigma p)
       =M(u,z)+\frac{p^2}{4}B_\kappa^2F(u,z).
 \label{eq:theta-symmetric-law}
\end{equation}
This is an algebraic cancellation, not a sign estimate.

\subsection{Convergent channels and physical recovery}
\label{sec:theta-recovery}

Let $h(u)=F(u,0)=e^u\Theta(e^{4u},0)$.  The normalized real and
imaginary period fields at $z=0$ are respectively
\[
 N_a(u,0)=\Phi_e(u)+\frac\pi4e^{-4u}(h'-h),\qquad
 N_b(u,0)=\Phi_e(u)-\frac\pi4e^{4u}(h'+h).
\]
Since $h(u)=e^u$ plus a superexponentially small remainder at positive
infinity and $h$ is even,
\[
 N_a(u,0)\sim-\frac\pi2e^{-5u}\quad(u\to-\infty),
 \qquad
 N_b(u,0)\sim-\frac\pi2e^{5u}\quad(u\to+\infty).
\]
Neither single normalized period field can be integrated over the
whole real line against $e^{tu^2+ixu}$ at positive real time.  The
growing carrier is not removed by the normalization in
\eqref{eq:theta-inserted-law}.  An analytic continuation or subtraction
would be a new operation requiring its own endpoint and domain ledger.

An exact convergent repair uses the symmetric short real period on
$u\ge0$ and the symmetric short imaginary period on $u<0$:
\begin{equation}
 \mathcal A(u,z)=
 \begin{cases}
 M(u,z)+e^{-4u}B_+^2F(u,z)/4,&u\ge0,\\
 M(u,z)-e^{4u}B_-^2F(u,z)/4,&u<0,
 \end{cases}
 \qquad B_\pm=\partial_z\mp\pi z.
 \label{eq:theta-short-field}
\end{equation}
For $u\ge0$, let
\[
 f(u,z)=F(u,z)-e^{u+\pi z^2/2},\qquad
 J_t(x,z)=\int_0^\infty e^{tu^2+ixu}f(u,z)\,du.
\]
All these half-line integrals converge locally uniformly with their
fixed derivatives. Define the full short-period transform by
\[
 \mathscr A_t(x,z)=\frac12\int_{\mathbb R}
       e^{tu^2+ixu}\mathcal A(u,z)\,du.
\]
It satisfies
\begin{equation}
 \mathscr A_t(x,z)=K_t(x,z)+\frac18
 \left[(B_w^2J_t)(x+4i,z)
              +(B_w^2J_t)(-x+4i,-iz)\right],
 \label{eq:theta-short-channel}
\end{equation}
where $B_w=\partial_w-\pi w$ acts on the own angular argument of $J$
before evaluation.  The imaginary spectral shifts supply the factor
$e^{-4u}$ on the half-line and are paid by the defining integrals.
In an unsymmetrized period relation, integration by parts additionally
produces the explicit lower endpoint $-f(0,z)$; it cancels in
\eqref{eq:theta-short-channel} under the exact opposite-period average.
The derivation, unsymmetrized endpoint formula, and complete compact-set
tail bounds are in \cite{ellipticnote}.

The convergent channel changes the observed source.  At $z=0$ put
\begin{equation}
 R(u)=2\pi^2e^u\sum_{n\ge1}n^2e^{-\pi n^2e^{4u}}>0,
 \qquad
 P_t(x)=\int_0^\infty e^{tu^2}R(u)\cos(xu)\,du.
 \label{eq:theta-correction-source}
\end{equation}
Then $\mathcal A(u,0)=\Phi_e(u)-R(|u|)$ and the exact physical
recovery is
\begin{equation}
 H_t=\mathscr A_t(\cdot,0)+P_t.
 \label{eq:theta-physical-recovery}
\end{equation}
For brevity write $\mathscr A_t(x)=\mathscr A_t(x,0)$.  Even a positive
$P_t$ does not control the physical Laguerre form, because
\begin{equation}
 \mathscr L_0(H)
 =\mathscr L_0(\mathscr A)+\mathscr L_0(P)
     +2\mathscr A'P'-\mathscr A P''-P\mathscr A''.
 \label{eq:theta-cross-balance}
\end{equation}
The derivatives and cross terms are part of the target.  There is no
sign estimate for them from positivity of $P$ alone.

The splice also explains a potentially misleading auxiliary
large-frequency term.  With $C(u)=-R(u)$ on $u\ge0$,
\[
 C'(0+)=\frac\pi4\bigl(16\Phi_e(0)+5h(0)\bigr)>0.
\]
Even reflection gives a cusp in the short-period source, and half-line
integration by parts yields, for fixed real $t$,
\[
 \mathscr A_t(x)=-\frac{C'(0+)}{x^2}+O_t(x^{-4}),
 \qquad
 P_t(x)=\frac{C'(0+)}{x^2}+O_t(x^{-4}).
\]
The original heated theta source is smooth and Schwartz, so its real
Fourier transform decays faster than every power.  The artificial
algebraic terms cancel in the physical recovery.  Their apparent
visibility cannot be transferred to $H$ separately.

\subsection{A positive Volterra correction still leaves both targets open}
\label{sec:theta-volterra}

Differentiating the source in
\eqref{eq:theta-correction-source} gives
\begin{equation}
 R'+5R=-4\pi e^{-4u}\Phi(u),\qquad
 R(u)=4\pi e^{-5u}\int_u^\infty e^v\Phi(v)\,dv.
 \label{eq:theta-volterra-law}
\end{equation}
The second formula uses the vanishing upper boundary $e^{5u}R(u)$.
The same Volterra definition can be made for any positive decaying
source $\psi$ on the positive half-line.  Writing
$r_{\psi,t}=e^{tu^2}R_\psi$ gives
\begin{equation}
 r_{\psi,t}''=e^{tu^2}\left(
   [(2tu-5)^2+2t]R_\psi
    +4\pi e^{-4u}[(9-4tu)\psi-\psi']\right).
 \label{eq:theta-correction-curvature}
\end{equation}
For $t\ge0$, the sufficient condition
$\psi'/\psi<9-4tu$ makes this strictly positive.  Rapid decay then
pays two integrations by parts, yielding
\begin{equation}
 P_{\psi,t}(x)=\frac1{x^2}\int_0^\infty
                r_{\psi,t}''(u)(1-\cos(xu))\,du>0
 \quad(x\ne0).
 \label{eq:theta-positive-correction}
\end{equation}
Positivity at $x=0$ follows from the source integral.  The genuine
theta summands give the logarithmic derivative bound
\[
 \frac{\Phi'}{\Phi}\le13-4\pi e^{4u},\qquad u\ge0,
\]
which meets this condition throughout $0\le t\le1$.

The same correction positivity survives the controls in
Theorem~\ref{thm:theta-finite-controls} when $t_0\in[0,1]$.  Indeed
their relative perturbations $q_X=\psi_X/\Phi_e-1$ and first
derivatives tend uniformly to zero.  For large $X$,
$|q_X'/(1+q_X)|\le1$, hence
\[
 \frac{\psi_X'}{\psi_X}\le14-4\pi e^{4u}<9-4tu
 \quad(0\le t\le1).
\]
The strict gap follows from $4u\le e^{4u}-1$:
\[
 9-4tu-(14-4\pi e^{4u})
 \ge(4\pi-1)e^{4u}-4>0.
\]
Thus these sources retain the scalar Volterra law, a positive
decreasing and convex heated correction source, and a positive Fourier
correction, while both physical stationary signs fail.  Those scalar
inputs cannot replace the missing arithmetic bound
\cite{localobstruction,volterranote}.

The remaining inequalities can be written without hiding a
normalization term.  Where $P>0$, set
\[
 Y=\frac HP=1+\frac{\mathscr A}P,\qquad
 p=\frac{P'}P,\qquad Z=Y'+pY=\frac{H'}P.
\]
Straight differentiation gives
\begin{equation}
 \frac{\mathscr L_0(H)}{P^2}=Y'^2-YY''-p'Y^2,
 \qquad
 \frac{\mathscr L_1(H)}{P^2}=Z'^2-ZZ''-p'Z^2.
 \label{eq:theta-normalized-targets}
\end{equation}
The two stationary equations become $Y'+pY=0$ and $Z'+pZ=0$.
Therefore proving the corresponding signs requires
\[
 Y'^2-YY''\ge p'Y^2\quad\hbox{when }Y'+pY=0,
 \qquad
 Z'^2-ZZ''\ge p'Z^2\quad\hbox{when }Z'+pZ=0.
\]
These are exact reformulations of the targets, not inequalities
already obtained from gluing or correction positivity.

\subsection{Why finite angular jets are not new independent anchors}
\label{sec:theta-angular-dictionary}

Termwise theta heat in scale and angle gives the balanced equation
\[
 F_u=\pi^{-1}F_{zz}-\pi z^2F.
\]
The same equation holds for $M$.  If $D_u=\partial_u|_z$, then
\[
 [z^{2m}]M=Q_m(D_u)\Phi_e,
 \qquad
 Q_{-1}=0,\quad Q_0=1,\quad
 Q_{m+1}=\frac{\pi D_uQ_m+\pi^2Q_{m-1}}{(2m+2)(2m+1)}.
\]
Full-line integration by parts, with the original smooth theta tails,
therefore gives
\begin{equation}
 [z^{2m}]K_t(x,z)=Q_m\bigl(i(2t\partial_x-x)\bigr)H_t(x).
 \label{eq:theta-angular-dictionary}
\end{equation}
Operator powers here retain the noncommutation of $x$ and
$\partial_x$.  In particular $Q_1=\pi D_u/2$, so
\[
 [z^2]K_t=\frac{i\pi}{2}(2tH_x-xH).
\]
At a joint zero this first angular coefficient vanishes automatically.
Higher coefficients need not vanish: they encode higher physical
derivatives.  The entire finite hierarchy is a scalar derivative
dictionary, rather than a separate nonvanishing source quantity.
Changing symmetric integer coefficients leaves the local balanced
equation and the real period intact, whereas the complete imaginary
law requires the recurrence in
\eqref{eq:theta-coefficient-defect}.  This is why the full gluing law
contains information beyond a finite jet list, while uniqueness alone
still supplies no signed physical inequality \cite{ellipticnote}.


\section{Corrected heat and operator machinery}
\label{sec:corrected-heat}

The earlier continuum, quantum, and Dirac notes
\cite{baker2017,baker2018heat,baker2018dirac} motivate the present
operator program.  Their exact semigroup identities survive, but several
operations must be distinguished.  If $\mathsf D$ is selfadjoint, spectral
calculus gives the state-space identity
\begin{equation}\label{eq:dirac-gaussian-average}
 e^{-a\mathsf D^2}
 =\frac1{\sqrt{4\pi a}}\int_{\mathbb R}
       e^{-s^2/(4a)}e^{is\mathsf D}\,ds,\qquad a>0,
\end{equation}
with a strong operator integral.  Averaging conjugations of a bounded
observable instead gives
\begin{equation}\label{eq:observable-dirac-average}
 \frac1{\sqrt{4\pi a}}\int_{\mathbb R}e^{-s^2/(4a)}
 e^{is\mathsf D}Be^{-is\mathsf D}\,ds
 =e^{-a\operatorname{ad}_{\mathsf D}^2}(B).
\end{equation}
The right side denotes the resulting observable semigroup; a generator
formula applies on its commutator domain.  A single first-order
propagator, a nonunitary similarity, state-space heat, and this unital
completely positive observable channel are different constructions.
None directly changes the sign of the backward Newman heat equation.

For example, suppose a balanced finite spin representation satisfies
$D^2=LI$ and $\operatorname{tr}_{\mathrm{spin}}D=0$, with normalized
spin trace and scalar $L$.  Its entire series is
\begin{equation}\label{eq:spin-trace-obstruction}
 \operatorname{tr}_{\mathrm{spin}}e^{sD}
 =\sum_{m\ge0}\frac{s^{2m}L^m}{(2m)!}
 =\cosh(s\sqrt L).
\end{equation}
For $L=-|k|^2$ this is $\cos(s|k|)$.  The coefficient of $s^2L$
is $1/2$, rather than the coefficient $1$ in $e^{s^2L}$.
Replacing $s$ by $\sqrt2s$ repairs that coefficient but gives $1/6$
instead of $1/2$ at $s^4L^2$.  More generally a finite positive spin
trace of $e^{iskS}$, $S=S^*$, is a finite characteristic function.
Its continuation at $k=iR$ has at most exponential growth in $|R|$;
it cannot equal the Gaussian $e^{-ak^2}$, whose continuation grows
as $e^{aR^2}$.  Gaussian averaging in
\eqref{eq:dirac-gaussian-average} is the valid replacement
\cite{smoothingnote,seedreview}.

The proposed constrained square in the Dirac source also needs its full
expansion.  If $\mathcal G$ is a density-linear, number-preserving field
and $c$ an annihilation field, with $\nu\ge0$, then
\[
 (\mathcal G+\sqrt\nu c)^*(\mathcal G+\sqrt\nu c)
 =\mathcal G^*\mathcal G
  +\sqrt\nu(\mathcal G^*c+c^*\mathcal G)+\nu c^*c.
\]
Its cross terms change particle number; on a fixed-number state their
expectations vanish instead of becoming a kinetic contribution.
Moreover $\mathcal G^*\mathcal G$ retains the density self-contraction
before it can be identified with a pair interaction.  These are
specific missing terms in that source's square completion, rather than
an exclusion of all Dirac constructions \cite{seedreview}.

\subsection{Translation channels and their product defect}
Let $U_y$, $y\in\mathbb R^d$, be a strongly continuous unitary
representation on a Hilbert space, and set
\begin{equation}\label{eq:gaussian-channel}
 \alpha_y(B)=U_yBU_y^*,\qquad
 P_a(B)=\int_{\mathbb R^d}K_a(y)\alpha_y(B)\,dy,\qquad
 K_a(y)=(4\pi a)^{-d/2}e^{-|y|^2/(4a)}.
\end{equation}
For bounded observables the integral is weak or ultraweak.  It is a
norm integral on the algebra with norm-continuous translation orbits.
Each $P_a$ is normal, unital, and completely positive, and
$P_aP_b=P_{a+b}$.  On the norm $C^2$ translation algebra put
\begin{equation}
 \delta_\ell B=\left.\partial_{y_\ell}\alpha_y(B)\right|_{y=0},
 \quad \mathcal L=\sum_\ell\delta_\ell^2,
 \quad \Gamma(B,C)=\sum_\ell(\delta_\ell B)^*\delta_\ell C.
\end{equation}
The product rule implies the exact defect
\begin{equation}\label{eq:channel-product-defect}
 P_a(B^*C)-P_a(B)^*P_a(C)
 =2\int_0^aP_{a-s}\bigl(\Gamma(P_sB,P_sC)\bigr)\,ds.
\end{equation}
Indeed, the derivative of
$P_{a-s}((P_sB)^*P_sC)$ is minus the integrand in
\eqref{eq:channel-product-defect}.  Its endpoint values give the
identity.  No smoothness is needed for the alternative formula
\begin{equation}\label{eq:channel-double-square}
 \begin{split}
 P_a(B^*C)-P_a(B)^*P_a(C)
 =\frac12\iint K_a(y)K_a(z)
 [\alpha_y(B)-\alpha_z(B)]^*
 [\alpha_y(C)-\alpha_z(C)]\,dy\,dz .
 \end{split}
\end{equation}
Expanding the four products proves it.  For $B=C$ it is a positive
operator, in agreement with the usual Schwarz inequality for unital
completely positive maps \cite{choi}.

The equality case also matters.  The defect for $B=C$ is
\[
 \int K_a(y)[\alpha_y(B)-P_a(B)]^*
                   [\alpha_y(B)-P_a(B)]\,dy.
\]
If it vanishes, testing on any vector makes a nonnegative continuous
scalar integrand vanish everywhere, since $K_a>0$.  Thus
$\alpha_y(B)=P_a(B)$ for all $y$, and $y=0$ gives $B=P_a(B)$.
Conversely translation invariance makes the defect zero.  Hence the
multiplicative domain is exactly the translation-invariant algebra
$\{U_y:y\in\mathbb R^d\}'$.  A single invertible similarity preserves
all products and cannot implement this channel on an algebra containing
a noninvariant observable together with its products.  For commuting
momenta $U_y=e^{iy\cdot P}$, the generator is
$-\sum_\ell[P_\ell,[P_\ell,B]]$, a square of derivations rather than
left multiplication by $-\sum_\ell P_\ell^2$.

\subsection{Independent density smoothing and normal ordering}
On $\bigwedge^N L^2(\mathbb R^d)$, with spin labels allowed, define
the bounded multiplication observables, for $a>0$,
\begin{equation}
 \rho_a(x)=\sum_{i=1}^NK_a(x-X_i),\qquad
 :\rho_a(x)^2:=\rho_a(x)^2-\sum_i K_a(x-X_i)^2.
\end{equation}
Normal ordering here removes exactly equal particle indices.  Gaussian
convolution gives
\begin{equation}\label{eq:normal-pair-Q}
 Q_a:=\int:\rho_a(x)^2:\,dx
 =\sum_{i\ne j}K_{2a}(X_i-X_j),\qquad
 \int\rho_a(x)^2\,dx=Q_a+NK_{2a}(0)I,
\end{equation}
so $0\le Q_a\le N(N-1)K_{2a}(0)I$.  The spatial integrals are
positive weak or quadratic-form integrals specified by these
multiplication formulas, rather than norm Bochner integrals over $x$.

Common translation of all particles yields
$P_s(\rho_a(x))=\rho_{a+s}(x)$, up to the immaterial translation
sign convention.  Smoothing the two density factors independently
and subtracting their same-index contraction produces $Q_{a+s}$.
Common translation of their already integrated pair instead gives
\begin{equation}\label{eq:common-translation-pair}
 P_s(Q_a)=Q_a,
\end{equation}
because it leaves $X_i-X_j$ unchanged.  The independent operation
diffuses the relative displacement; the common operation does not.
Spatially integrating the ordinary Schwarz defect gives only
\begin{equation}\label{eq:normal-order-nonmonotonicity}
 Q_a-Q_{a+s}+N[K_{2a}(0)-K_{2(a+s)}(0)]I\ge0.
\end{equation}
Deleting the positive self-contraction does not imply $Q_a\ge Q_{a+s}$.
For two distinct particle positions separated by $r>0$ the pair value is
\begin{equation}
 Q_a=2(8\pi a)^{-d/2}e^{-r^2/(8a)},\qquad
 \partial_a\log Q_a=-\frac d{2a}+\frac{r^2}{8a^2},
\end{equation}
which increases for $a<r^2/(4d)$ and decreases afterward.  This also
gives form counterexamples using wavefunctions supported near such
separations.  In $d=3$ the positive integrated Coulomb cutoff is a
different quantity:
\begin{equation}\label{eq:coulomb-scale-cutoff}
 2\pi e^2\int_b^\infty Q_{\tau/2}\,d\tau
 =e^2\sum_{i<j}
 \frac{\operatorname{erf}(|X_i-X_j|/(2\sqrt b))}{|X_i-X_j|},
 \qquad b>0.
\end{equation}
It decreases with $b$.  Its monotonicity supplies no monotonicity of
the scale density $Q_a$.  Unsmeared self and pair terms cannot be
split into individually divergent integrals and then discarded.
The bounded fermionic examples and verification details are recorded
in \cite{smoothingnote}; the formulas above are the relevant mechanism.

\section{Residual-corrected integrated quantum paths}
\label{sec:quantum-paths}
This secondary application tests the same discipline of retaining a
discarded term.  It does not imply a stationary theta sign.  In the
2018 quantum construction \cite{baker2018heat}, if
$\mathsf A=\mathsf A^*\ge\alpha I>0$, $\rho$ is any Hilbert-space
vector, and $u(s)=e^{-s\mathsf A}\rho$, then
\begin{equation}\label{eq:infinite-quantum-identities}
 \Psi=\int_0^\infty u(s)\,ds=\mathsf A^{-1}\rho,\qquad
 \langle\Psi,\mathsf A\Psi\rangle
 =2\int_0^\infty\|u(s)\|^2\,ds,\qquad
 \|\Psi\|^2=4\int_0^\infty s\|u(s)\|^2\,ds.
\end{equation}
Spectral calculus proves these equalities by integrating
$e^{-2s\lambda}$ and $s e^{-2s\lambda}$ for $\lambda>0$.
The gap ensures convergence and inverse domains.  An integral of
Slater determinants can define a correlated trial state even though
each slice is a determinant.  Replacing its evolution without a
correction, however, loses the diagonal quotient's variational meaning.

\begin{theorem}[Finite-time path identity]\label{thm:residual-path}
Let $\mathsf A$ be selfadjoint and let $u:[0,T]\to\operatorname{Dom}
\mathsf A$ be continuously differentiable in Hilbert norm and
continuous in the graph norm.  Set
\begin{equation}\label{eq:path-defects}
 r(s)=u'(s)+\mathsf A u(s),\quad
 U(s)=\int_s^T u(v)\,dv,\quad
 d(s)=\int_s^T r(v)\,dv-u(T),\quad \Psi=U(0).
\end{equation}
Then $\mathsf A U=u+d$, and
\begin{align}
 E:=\langle\Psi,\mathsf A\Psi\rangle
 &=2\int_0^T\|u\|^2\,ds
   +2\operatorname{Re}\int_0^T\langle u,d\rangle\,ds,
   \label{eq:residual-energy}\\
 N:=\|\Psi\|^2
 &=4\int_0^T s\|u\|^2\,ds
   +4\operatorname{Re}\int_0^T s\langle u,d\rangle\,ds
   -2\operatorname{Re}\int_0^T\langle U,d\rangle\,ds.
   \label{eq:residual-norm}
\end{align}
In finite dimension it suffices that $u$ be continuously differentiable.
\end{theorem}
\begin{proof}
Integration of $r$ gives $\mathsf A U=u+d$.  The real function
$F(s)=\langle U(s),\mathsf A U(s)\rangle$ satisfies
\[
 F'=-2\operatorname{Re}\langle u,\mathsf A U\rangle
    =-2\|u\|^2-2\operatorname{Re}\langle u,d\rangle,
 \qquad F(T)=0,
\]
which proves \eqref{eq:residual-energy}.  Since $U'=-u$,
\[
 N=2\operatorname{Re}\int_0^T\langle U,u\rangle\,ds
  =2\int_0^T F(s)\,ds
   -2\operatorname{Re}\int_0^T\langle U,d\rangle\,ds.
\]
The identity $\int_0^TF=-\int_0^TsF'$ proves
\eqref{eq:residual-norm}.  The graph-norm hypothesis justifies
integration in $\operatorname{Dom}\mathsf A$ and all derivatives.
\end{proof}

For $\Psi\ne0$, $E/N$ is its exact Rayleigh quotient regardless of
the quality of the path; hence $\inf\sigma(\mathsf A)\le E/N$.
For a quantum Hamiltonian $\mathsf H$ and a declared scalar energy shift
$E_{\mathrm{sh}}$, set $\mathsf A=\mathsf H+E_{\mathrm{sh}}I$ and subtract
$E_{\mathrm{sh}}$ from the quotient.  Even when $r=0$, finite time
leaves $d=-u(T)$, so omitting the endpoint changes both the energy
and norm.  For the scalar example $\mathsf A=aI$ and approximate
$u(s)=e^{-bs}v$, $a,b>0$, the uncorrected infinite-time diagonal
quotient is $b$ whereas the Rayleigh quotient is $a$.  Thus $b<a$
violates the claimed upper bound.  These elementary corrections and
their finite-time formulation are documented in \cite{residualnote}.

Write $Q_0=\int\|u\|^2$, $Q_1=\int s\|u\|^2$,
$D_0=\int\|d\|^2$, $D_1=\int s\|d\|^2$, and
$V_0=\int\|U\|^2$, all on $[0,T]$.  Cauchy--Schwarz gives
\begin{equation}\label{eq:generic-residual-bounds}
 |E-2Q_0|\le\delta_E:=2\sqrt{Q_0D_0},\qquad
 |N-4Q_1|\le\delta_N:=4\sqrt{Q_1D_1}+2\sqrt{V_0D_0}.
\end{equation}
If $\mathsf A\ge0$ and $4Q_1>\delta_N$, then
\begin{equation}\label{eq:generic-rayleigh-upper}
 \inf\sigma(\mathsf A)\le\frac EN
 \le\frac{2Q_0+\delta_E}{4Q_1-\delta_N}.
\end{equation}
The strict denominator condition is a substantive gate.  Bounding
the correction terms separately loses their cancellations; the path
identity alone does not make this lower bound positive.

\subsection{A positive-interaction estimate}
In finite dimension let $\mathsf B=\mathsf B^*\ge\beta I>0$,
$\mathsf A=\mathsf B+\mathsf V$, and $\mathsf V\ge0$ with
$\mathsf V^2\ge\nu\mathsf V$, $\nu>0$.  Use the declared fixed
surrogate path
\begin{equation}
 u(s)=e^{-s\mathsf B}\rho,\qquad
 r(s)=\mathsf V u(s),\qquad
 R=\int_0^Tr(s)\,ds=\mathsf V\Psi.
\end{equation}
Let $\ell=\sqrt{4Q_1}-\|u(T)\|/\beta$.  If $\ell>0$, then
\begin{equation}\label{eq:positive-interaction-upper}
 \inf\sigma(\mathsf A)\le\frac{\langle\Psi,\mathsf A\Psi\rangle}
 {\|\Psi\|^2}
 \le\frac{2Q_0+\|R\|^2/\nu}{\ell^2}.
\end{equation}
To prove the denominator bound, put $S=\mathsf B^{-1}\rho$.
Spectral calculus gives
$\|S\|^2=4\int_0^\infty s\|u(s)\|^2ds\ge4Q_1$ and
$\Psi=S-\mathsf B^{-1}u(T)$, so $\|\Psi\|\ge\ell$.
For the numerator,
\begin{equation}
 E_{\mathsf B}:=\langle\Psi,\mathsf B\Psi\rangle
 =2Q_0-2\operatorname{Re}\langle\Psi,u(T)\rangle\le2Q_0.
\end{equation}
The endpoint overlap is nonnegative because its spectral weight is
$(1-e^{-T\lambda})e^{-T\lambda}/\lambda$.  Finally
$\langle\Psi,\mathsf V\Psi\rangle\le\|\mathsf V\Psi\|^2/\nu$
proves \eqref{eq:positive-interaction-upper}.  No commutation between
$\mathsf B$ and $\mathsf V$ is required.  The exact endpoint relation
is $\mathsf A\Psi=\rho+R-u(T)$.  A signed interaction or a positive
interaction with spectrum accumulating at zero does not meet the gap
hypothesis.  When $\mathsf V=\nu P$ for a projection, the interaction
estimate is an equality \cite{hubbardnote}.

The interaction gap does not eliminate the main representation
bottleneck:
\begin{equation}\label{eq:residual-cross-kernel}
 \|R\|^2=\int_0^T\!\int_0^T
 \langle\mathsf V u(s),\mathsf V u(v)\rangle\,ds\,dv.
\end{equation}
Residual cross overlaps remain, and $\mathsf V$ need not preserve
a single determinant.  Replacing this norm of an integral by the
integral of norms gives a valid but generally weaker estimate.  A
useful quantum application requires a representation in which these
corrections are economical; an exact residual reformulation alone
does not establish an advantage over direct Rayleigh evaluation.
The bounded comparison and its limitations are in \cite{hubbardnote}.

\subsection{Elementary diagonal deductions and possible alternatives}
For the same fixed selfadjoint $\mathsf B$, put $q(s)=\|u(s)\|^2$.
The semigroup identity
$\langle u(s),u(v)\rangle=q((s+v)/2)$ gives the exact finite-time
formulas
\begin{align}
 N&=4\int_0^{T/2}s q(s)\,ds
       +4\int_{T/2}^T(T-s)q(s)\,ds,
       \label{eq:triangular-diagonal-norm}\\
 E_{\mathsf B}&=2\int_0^{T/2}q(s)\,ds
                  -2\int_{T/2}^Tq(s)\,ds.
       \label{eq:triangular-diagonal-energy}
\end{align}
For \eqref{eq:triangular-diagonal-norm}, integrate the double overlap
over $[0,T]^2$; the level set $s+v=r$ has measure $r$ for $r<T$
and $2T-r$ for $r>T$.  Substitute $r=2s$.
For \eqref{eq:triangular-diagonal-energy}, use
$\langle u(s),\mathsf B u(v)\rangle=-q'((s+v)/2)/2$ and integrate
the same triangular weights by parts.  Spectral calculus justifies
the result for arbitrary $\rho$; a Hilbert-norm derivative at $s=0$
is unnecessary.  These are elementary deductions of fixed-generator
evolution, rather than a novelty claim.  In particular the exact
denominator is obtainable from diagonal data, with strictly positive
weights when $\rho\ne0$.  This observation does not apply to an
arbitrary time-dependent approximate path.

A possible gap-free replacement for the interaction bound is
\begin{equation}\label{eq:gapfree-diagonal-interaction}
 \langle\Psi,\mathsf V\Psi\rangle
 =\left\|\int_0^T\mathsf V^{1/2}u(s)\,ds\right\|^2
 \le\left(\int_0^T
       \sqrt{\langle u(s),\mathsf V u(s)\rangle}\,ds\right)^2.
\end{equation}
This is the triangle inequality and requires only $\mathsf V\ge0$.
For an unbounded interaction it is a form identity under integrability
in the $\mathsf V^{1/2}$ graph norm.  A form bound
$\mathsf V\le a\mathsf B+bI$, $a,b\ge0$, further bounds the
integrand by
$\sqrt{a\langle u,\mathsf B u\rangle+bq}$.
Combining \eqref{eq:triangular-diagonal-norm} with
\eqref{eq:gapfree-diagonal-interaction} could replace a two-time
interaction kernel by diagonal interaction information.  Its sharpness
and useful applications remain to be established; it is a candidate
estimate, not a completed many-body method.

\subsection{An analytic benchmark model}
For context, the benchmark in \cite{hubbardnote} uses the two-site
Hubbard model with hopping $1$, repulsion $4$, and
$N_\uparrow=N_\downarrow=1$.  In the basis
$(L_\uparrow L_\downarrow,L_\uparrow R_\downarrow,
R_\uparrow L_\downarrow,R_\uparrow R_\downarrow)$ it has
\begin{equation}
 \mathsf H=\mathsf H_{\mathrm{hop}}+\mathsf V,\quad
 \mathsf H_{\mathrm{hop}}=\begin{pmatrix}0&-1&-1&0\\-1&0&0&-1\\-1&0&0&-1\\0&-1&-1&0\end{pmatrix},
 \qquad \mathsf V=\operatorname{diag}(4,0,0,4).
\end{equation}
The symmetric two-dimensional block is
$\left(\begin{smallmatrix}4&-2\\-2&0\end{smallmatrix}\right)$,
giving the exact ground energy $2-2\sqrt2$.  Every normalized
sector-preserving single Slater determinant is a product of an up
orbital and a down orbital.  Let their site polarizations be $z,w$.
Minimizing their phases gives
\begin{equation}
 E_{\mathrm{Slater}}(z,w)
 =-\sqrt{1-z^2}-\sqrt{1-w^2}+2(1+zw).
\end{equation}
Minimization over $w$ and $s=\sqrt{1-z^2}\in[0,1]$ yields
$E_{\mathrm{Slater}}\ge2-s-\sqrt{5-4s^2}\ge-1/2$, because
\[
 (5/2-s)^2-(5-4s^2)=5(s-1/2)^2\ge0.
\]
Equality occurs at $s=1/2$ with opposite polarizations.  The trial
family integrates the determinant slices
$e^{-s(\mathsf H_{\mathrm{hop}}+3I)}(a_\uparrow\otimes b_\downarrow)$, and its integral
can have Schmidt rank two.  This supplies an analytic comparison
class and a nonzero correlated-energy target.  The actual certified
trial energies, residual enclosures, and evaluation comparison are
reported only in the benchmark note \cite{hubbardnote}; no numerical
table is needed for the operator conclusions here.

\section{Bottlenecks and open construction targets}\label{sec:bottlenecks}
The formulas distinguish several difficulties that require different
remedies. This section records possible directions without selecting a
new implementation or assuming a missing positivity statement.

\subsection{From a generator square to a stationary observation}
Equations \eqref{eq:local-square-defect}--\eqref{eq:integrated-energy}
show the basic loss of information. The squared first-order derivative
controls an integrated energy, while the stationary Laguerre form is
spatial curvature of $h_j^2$. The paired representation makes the same
loss visible as the negative part of an oscillatory spectral multiplier.
The right target remains the unnormalized physical observation.
Changing its normalization requires the correction in
\eqref{eq:normalized-critical-energy}.

A first possible direction is a comparison law for critical values.
Equation \eqref{eq:critical-higher-jets} identifies the exact higher-jet
combination requiring an arithmetic bound. Merely assuming appropriate
critical curvature, real-root interlacing, or log-concavity of the physical
transform would import the desired conclusion. A useful comparison lemma
must come from the complete theta preparation and remain valid near the
candidate degeneracies. The weaker stationary-loss target
\eqref{eq:bounded-loss-target} offers a second endpoint for such a lemma:
it permits an adverse sign provided the loss is bounded relative to the
vanishing joint first jet, locally through every possible collision.

\subsection{Why a nonlocal square alone does not repair the core identity}
The local multiplier obstruction produces a negative compact-support
quadratic form of the same $K_x-P_x$. Consequently this operator cannot
be equal to \emph{any} positive quadratic form on that full core, even if
its proposed square factor is nonlocal. A viable construction must retain
an additional constraint term that is nonzero on general vectors, or
prove positivity on a precisely declared arithmetic family.

The distinction is important for gluing. A gluing term may vanish on the
actual theta state without being a zero operator. Its one-state vanishing
cannot justify removing it from a full-domain identity. Similarly,
compressing the target to a state family already uniquely selected by
full gluing does not prove positivity of that compression. The latter is
still the missing signed estimate unless an independent factorization or
domination has been constructed.

\subsection{A possible adjoint gluing certificate}
An explicit lifted identity could have the shape
\begin{equation}\label{eq:open-adjoint-gluing}
\begin{split}
 \mathscr L_j(H;t,x)={}&\|\mathcal Q_{j,t,x}\Psi_t\|^2
 +2\Rea\langle\mathcal G_t\Psi_t,
                  \mathcal R_{j,t,x}\Psi_t\rangle\\
 &+H^{(j+1)}(t,x)M_{j,t,x}+E_{j,t,x}.
\end{split}
\end{equation}
This is a research target, not an identity constructed in this paper.
The lifted state, its domain, and its physical readout must be explicit.
The operator $\mathcal G_t$ must encode the full imaginary Jacobi law
and its transported insertion, rather than finitely many scalar tests.
For genuine theta $\mathcal G_t\Psi_t=0$; the relevant stationary
constraint eliminates the next term. The square must then dominate every
negative part of $E_j$, or $E_j$ must satisfy an independently proved
bound sufficient for \eqref{eq:bounded-loss-target}.

The coefficient recurrence $c_n-c_{n+1}=0$ suggests a discrete adjoint
summation or telescoping identity. The substantive intermediate lemma
would construct $\mathcal R_j$ from that summation and prove convergence
of the infinite coefficient pairing after the physical insertion. It must
retain the moving-period connection, zero-mode subtraction, endpoint
partners and both source tails. A finite symbolic operator can encode the
infinite recurrence: the finite-observation obstruction does not rule out
such a construction. A square obtained by taking the unknown square root
of $\mathscr L_j$ or dividing by $H$ at unlocated zeros would not supply
an independent certificate.

A simple Gaussian embedding illustrates both the availability and the
limitation of a gluing square. For fixed $r>0$ and coefficient sequences whose Gaussian embedding
belongs to $\ell^2(\mathbb Z)$, set
\begin{equation}
 g_n=c_n e^{-\pi rn^2/2},\qquad
 (Q_rf)_n=f_{n+1}-e^{-\pi r(n+1/2)}f_n.
 \label{eq:gluing-weighted-shift}
\end{equation}
On $\ell^2(\mathbb Z)$, $Q_r$ is a bounded bilateral shift minus a
closed multiplication operator, on the domain of that multiplier.
It is closed and densely defined. For $c_n=1$ the Gaussian vector lies
in this domain and satisfies $Q_rg=0$, so $Q_r^*Q_r$ supplies a positive
quadratic form with the arithmetic preparation as a null vector.
But this is a square of the preparation constraint, not the physical
Laguerre observation. Its coefficients grow on one index tail, and the
radial dependence $r=e^{4u}$ introduces additional derivatives when the
physical insertion is applied. An annihilating Hamiltonian by itself
therefore does not provide the missing signed physical projection.

\subsection{A possible transport between the signed sectors}
For $j=0$ let $A$ multiply by $u-v$, let $S$ multiply by $u+v$, and
write
\begin{equation}
 C_\pm=\frac{|\cos(xS)|\pm\cos(xS)}2,\qquad
 p=C_+^{1/2}AW_t,\quad n=C_-^{1/2}AW_t.
\end{equation}
Then $2\mathscr L_0=\|p\|^2-\|n\|^2$. Another possible intermediate
lemma would prescribe, from gluing, a map $T_{t,x}$ with
\begin{equation}\label{eq:open-contraction}
 n=T_{t,x}p+\text{explicit stationary and gluing defects},
 \qquad \|T_{t,x}\|\le1
\end{equation}
on a declared domain. At genuine stationarity this proves domination.
For $j=1$ the signed split must instead use $(-uv)\cos[x(u+v)]$;
the first-level cosine split cannot simply be reused.

Constructing a rank-one map from the already given vectors $p,n$ has
norm $\|n\|/\|p\|$ and merely repackages the target. The transport must
be prescribed independently and its norm proved before the physical sign
is inferred. Equation \eqref{eq:gluing-weighted-shift} shows why an exact
lattice shift is not automatically a contraction under the Gaussian and
physical heat weights. The admissible norm and all weight ratios are part
of the missing lemma.

\subsection{Residual cancellation and the interaction class}
The quantum path identities are exact. Their generic upper estimate loses
usefulness when signed cross terms are separately replaced by absolute
errors, especially when the lower norm bound ceases to be positive.
Equations \eqref{eq:triangular-diagonal-norm} and
\eqref{eq:triangular-diagonal-energy} remove that denominator loss within
fixed selfadjoint surrogate evolution, but do not apply automatically to
an arbitrary time-dependent approximation.

The positive-interaction residual bound uses a known gap on the positive
spectrum of the interaction. Its residual norm retains the two-time cross
kernel in \eqref{eq:residual-cross-kernel}; the rewriting does not remove
the overlap problem. The gap-free diagonal inequality
\eqref{eq:gapfree-diagonal-interaction}, a sufficiently sharp relative
form bound, or a block Gram estimate preserving signed correlations are
possible alternatives. Their utility depends on obtaining those inputs
without reintroducing more expensive direct energy and norm calculations.
The corresponding theorem would concern a declared quantum approximation,
not the theta stationary sign.

\subsection{Normal ordering and coverage remain separate obligations}
The positive Gaussian channel defect controls an ordinary operator square.
After pair-index subtraction, only
\eqref{eq:normal-order-nonmonotonicity} follows. A first-order construction
for normal-ordered interactions must therefore retain and control the
self contraction, and distinguish independent relative diffusion from
common translation. The corrected machinery settles the operator
implementation issue; it supplies no positivity transfer to Newman time.

For the theta application, a signed estimate must also cover the required
positive-time candidate set and its predecessor neighborhoods. Source
cutoffs and normalization must remain compatible across chart boundaries.
A theorem covering every possible candidate is sufficient even if it
leaves unrelated regions untreated; that candidate coverage itself needs
a proof. No local favorable sign, fixed positive time, or finite sample
provides this endpoint coverage. The global conditional implications are
precise, while their required genuine-theta estimate remains open.

\section{Supporting notes and computational records}\label{sec:records}
The analytic formulas above are the subject of this manuscript. The
supporting notes retain the bounded computations and their full error
ledgers. The stationary theta and Euler-control note
\cite{stationarycontrols} proves local signs on declared stationary curves
using complete theta, shifted-source and parameter remainders. Those
results test the physical readout and distinguish the control, but do not
supply global collision coverage.
The Hubbard note \cite{hubbardnote} records the selected variational
comparison, residual enclosure and matched work accounting for the
declared fixed-number model. Its useful bounded energy certificate and
its representation-dependent cost limitation are not a scalable
performance theorem. The smoothing note \cite{smoothingnote} retains the
exact finite CAR checks. No numerical table or measured margin is used
as a premise of the analytic theorems in this manuscript.

\section*{Acknowledgements and preparation record}
This draft was prepared for Edward Baker with substantial LLM assistance
using OpenAI Codex, recorded model \texttt{gpt-6.1-sol}, reasoning effort
\textbf{ultra}, on 14 October 2026. Parallel LLM drafting and mathematical
checks assisted organization, derivations, references and source review.
Earlier supporting notes retain their own model and effort provenance.
The author's 2017--2018 source notes remain his work and are cited
separately. Internal LLM checks are not independent specialist validation.
The conditional theta construction remains open, and no literature
priority claim is assigned to the elementary heat, spectral-calculus or
semigroup identities collected here.


\begin{thebibliography}{99}
\bibitem{polymath}
D. H. J. Polymath,
\emph{Effective approximation of heat flow evolution of the Riemann
$\xi$ function, and a new upper bound for the de Bruijn--Newman constant},
Research in the Mathematical Sciences \textbf{6} (2019), article 31.
\url{https://arxiv.org/abs/1904.12438v2}.
The normalization, threshold framework, Theorem 1.5(i), and Proposition 3.1
are the imported inputs used here.

\bibitem{newmanwu}
C. M. Newman and W. Wu,
\emph{Constants of de Bruijn--Newman type in analytic number theory and
statistical physics}, Bulletin of the American Mathematical Society
\textbf{57} (2020), 595--614.
\url{https://arxiv.org/abs/1901.06596}.

\bibitem{rodgerstao}
B. Rodgers and T. Tao,
\emph{The de Bruijn--Newman constant is non-negative},
Forum of Mathematics, Pi \textbf{8} (2020), e6.
\url{https://arxiv.org/abs/1801.05914}.

\bibitem{dlmf}
NIST Digital Library of Mathematical Functions,
Sections 20.2 and 20.7, especially equations 20.2.8 and 20.7.32,
Jacobi quasiperiodicity and inversion.
\url{https://dlmf.nist.gov/20.2.E8};
\url{https://dlmf.nist.gov/20.7.E32}.

\bibitem{choi}
M.-D. Choi,
\emph{A Schwarz inequality for positive linear maps on $C^*$-algebras},
Illinois Journal of Mathematics \textbf{18} (1974), 565--574.
\url{https://doi.org/10.1215/ijm/1256051007}.
The Gaussian channel identities in this manuscript have direct proofs;
this reference provides the general Schwarz-inequality context.

\bibitem{baker2017}
Edward B. Baker III,
\emph{Notes on the continuum sigma limit}, 22 November 2017,
unpublished author-supplied notes, source file
\texttt{Continuum\_sigma.pdf}. Source identity and original location are
recorded in the investigation's \nolinkurl{../SOURCE_DOCUMENTS.json}.

\bibitem{baker2018heat}
Edward B. Baker III,
\emph{Notes on the continuum sigma limit}, 28 February 2018,
unpublished author-supplied notes, source file
\texttt{Heat\_Quantum.pdf}. Source manifest as in \cite{baker2017}.

\bibitem{baker2018dirac}
Edward B. Baker III,
\emph{Notes on the continuum sigma limit}, 17 August 2018,
unpublished author-supplied notes, source file
\texttt{Dirac\_operator\_flow.pdf}. Source manifest as in \cite{baker2017}.

\bibitem{parentcollision}
\emph{Collision Exclusion and Collective Attraction in the Riemann Heat
Flow: Reductions, Quantitative Tools, and Endpoint Obstructions},
drafted for Edward Baker, 9 October 2026, project manuscript, Sections 1--2.
Repository path:
\nolinkurl{papers/quasi-rh-exponent-descent/newman_collisions/newman_collision_reductions.tex}.

\bibitem{upstreamtransport}
\emph{Enlarged-state transport and visibility for the Newman heat flow},
drafted for Edward Baker, 10 October 2026, project manuscript,
Sections 5--6 for the stationary criterion and paired Laguerre kernels.
Repository path:
\nolinkurl{papers/quasi-rh-exponent-descent/newman_collisions/17_enlarged_state_transport/enlarged_state_transport_and_heat_flow.tex}.

\bibitem{stationarynote}
\emph{A stationary-sign criterion and complete theta targets},
10 October 2026, upstream project note.
Repository path:
\nolinkurl{papers/quasi-rh-exponent-descent/newman_collisions/17_enlarged_state_transport/notes/5_STATIONARY_SIGN_CRITERION_AND_THETA_TARGETS_20261010.md}.

\bibitem{kernelnote}
\emph{Complete Laguerre kernels and the triple-zero limit},
10 October 2026, upstream project note.
Repository path:
\nolinkurl{papers/quasi-rh-exponent-descent/newman_collisions/17_enlarged_state_transport/notes/7_COMPLETE_LAGUERRE_KERNELS_AND_TRIPLE_ZERO_LIMIT_20261010.md}.

\bibitem{boundedloss}
\emph{Next steps and a bounded loss criterion for heat collisions},
10 October 2026, upstream project note.
Repository path:
\nolinkurl{papers/quasi-rh-exponent-descent/newman_collisions/17_enlarged_state_transport/notes/16_NEXT_STEPS_AND_BOUNDED_LOSS_COLLISION_CRITERION_20261010.md}.

\bibitem{eulercontrols}
\emph{Euler-Factor Controls for the de Bruijn--Newman Constant:
Certified and Asymptotic Lower Bounds}, drafted for Edward Baker,
14 October 2026, investigation 19 manuscript.
Repository path:
\nolinkurl{papers/investigations-2/active/19_euler_factor_controls/manuscript/euler_factor_controls.tex}.

\bibitem{ellipticnote}
\emph{Elliptic insertion transport and its heat obstructions},
13 October 2026, investigation 02 research note.
Repository path:
\nolinkurl{papers/investigations-2/active/02_angular_jacobi_gluing/notes/ELLIPTIC_INSERTION_TRANSPORT_AND_HEAT_OBSTRUCTIONS_20261013.md}.

\bibitem{volterranote}
\emph{Complete-channel Volterra reduction and stationary-sign controls},
13 October 2026, investigation 02 research note.
Repository path:
\nolinkurl{papers/investigations-2/active/02_angular_jacobi_gluing/notes/3_COMPLETE_CHANNEL_VOLTERRA_CONTROLS_20261013.md}.

\bibitem{seedreview}
\emph{Review of the three heat and operator source notes},
14 October 2026, investigation 20 source review.
\nolinkurl{../reviews/THREE_SOURCE_NOTES_REVIEW_20261014.md}.

\bibitem{diracnote}
\emph{Dirac squares and stationary Laguerre energy},
14 October 2026, investigation 20 research note.
\nolinkurl{../notes/DIRAC_SQUARES_AND_STATIONARY_LAGUERRE_ENERGY_20261014.md}.

\bibitem{residualnote}
\emph{Residual corrected heat trajectory identities},
14 October 2026, investigation 20 research note.
\nolinkurl{../notes/RESIDUAL_CORRECTED_HEAT_TRAJECTORY_IDENTITIES_20261014.md}.

\bibitem{localobstruction}
\emph{Local multipliers and obstructions from finite observations},
14 October 2026, investigation 20 research note.
\nolinkurl{../notes/LOCAL_MULTIPLIER_AND_FINITE_OBSERVATION_OBSTRUCTIONS_20261014.md}.

\bibitem{smoothingnote}
\emph{Gaussian translation defects independent smoothing and normal ordering},
14 October 2026, investigation 20 research note.
\nolinkurl{../notes/GAUSSIAN_TRANSLATION_DEFECTS_AND_NORMAL_ORDERING_20261014.md}.

\bibitem{hubbardnote}
\emph{Certified Hubbard heat trial and its cost limitation},
14 October 2026, investigation 20 research note and bounded computational record.
\nolinkurl{../notes/HUBBARD_HEAT_RESIDUAL_BENCHMARK_20261014.md}.

\bibitem{stationarycontrols}
\emph{Certified stationary theta and Euler controls},
14 October 2026, investigation 20 research note and bounded computational record.
\nolinkurl{../notes/2_CERTIFIED_STATIONARY_THETA_AND_EULER_CONTROLS_20261014.md}.
\end{thebibliography}
\end{document}

